\documentclass[11pt,a4paper]{article}
\usepackage[margin=2.2cm]{geometry}
\usepackage[T1]{fontenc}
\usepackage[utf8]{inputenc}
\usepackage{lmodern}
\usepackage{amsmath,amssymb,mathtools}
\usepackage{hyperref}
\usepackage{enumitem}

\title{Quartic Rovibrational Distortion Project\\
\large Current Strategy and Project Status}
\author{}
\date{\today}

\begin{document}
\maketitle

\section{Project Goal}
The goal of the project is to derive and evaluate quartic centrifugal distortion constants from a Watson-type rovibrational Hamiltonian expanded in normal coordinates. The long-term target is a practical program that takes molecular data (ultimately from Gaussian or equivalent sources), computes the rovibrational couplings, and returns quartic distortion constants in a Watson representation.

At the formal level, the project aims to isolate the quartic rotational block of the effective Hamiltonian and express it in terms of:
\begin{itemize}[leftmargin=*]
\item first-order rovibrational inertia derivatives \( \mu_{\alpha\beta,k} \),
\item second-order rovibrational inertia derivatives \( \mu_{\alpha\beta,kl} \),
\item cubic force constants \( \Phi_{ijk} \),
\item quartic force constants \( \Phi_{ijkl} \),
\item harmonic frequencies \( \omega_i \).
\end{itemize}

\section{Current Strategy}
The work is currently split into two complementary tracks.

\subsection{Track A: Full BCH / Van Vleck engine}
The file
\begin{center}
\texttt{derive\_watson\_quartic\_vanvleck.py}
\end{center}
implements a general symbolic Lie-transform / Van Vleck workflow. It uses:
\begin{itemize}[leftmargin=*]
\item bosonic creation and annihilation operators,
\item automatic normal ordering,
\item non-commutative rotational operators,
\item order-by-order elimination of vibrationally off-diagonal terms,
\item projection of the final quartic rotational block onto a commuting quartic basis and then onto Watson constants.
\end{itemize}

This engine is valuable because it provides a formal reference and a benchmark. However, for the physically relevant case of several vibrational modes and fourth perturbative order, the general BCH expansion becomes computationally expensive.

\subsection{Track B: Direct channel decomposition}
To avoid the combinatorial cost of the full BCH engine, the project now also uses a cheaper and more physical decomposition implemented in
\begin{center}
\texttt{quartic\_channels.py}
\end{center}
This second module isolates individual perturbative channels contributing to the quartic rotational block:
\begin{itemize}[leftmargin=*]
\item \(H_{12}H_{12}\): harmonic rovibrational second-order block,
\item \(H_{22}\): quadratic inertia block,
\item \(H_{12}H_{30}\): mixed cubic--rovibrational block,
\item \(H_{30}H_{30}\): cubic--cubic anharmonic block,
\item \(H_{40}\): quartic force-field block.
\end{itemize}

The key strategic idea is:
\begin{quote}
instead of building the entire effective Hamiltonian with a general BCH algebra, derive and validate each physically relevant channel separately, then sum the channels at the level of the quartic tensor \( \tau \).
\end{quote}

\section{Why This Strategy}
The BCH engine is formally general but expensive. The direct-channel approach is more suitable for a real molecular workflow because:
\begin{itemize}[leftmargin=*]
\item it matches the physical partition of perturbative contributions,
\item it makes debugging easier,
\item it allows channel-by-channel validation against known limits,
\item it is closer to the final numerical program that will use molecular input data.
\end{itemize}

In practice, the BCH engine is now used primarily as a reference generator for small benchmark cases, while the channel module is becoming the main computational route.

\section{Current Formal Picture}
The project is organized around the following perturbative classes:
\begin{align}
H_{12} &\equiv \frac12 \sum_{\alpha\beta k} \mu_{\alpha\beta,k} q_k J_\alpha J_\beta,\\
H_{22} &\equiv \frac14 \sum_{\alpha\beta kl} \mu_{\alpha\beta,kl} q_k q_l J_\alpha J_\beta,\\
H_{30} &\equiv \frac16 \sum_{ijk} \Phi_{ijk} q_i q_j q_k,\\
H_{40} &\equiv \frac1{24} \sum_{ijkl} \Phi_{ijkl} q_i q_j q_k q_l.
\end{align}

The quartic rotational block is then decomposed into additive contributions of the form:
\begin{equation}
\tau^{(4)} = \tau^{(H_{12}H_{12})} + \tau^{(H_{22})} + \tau^{(H_{12}H_{30})} + \tau^{(H_{30}H_{30})} + \tau^{(H_{40})}.
\end{equation}

This is the object that is eventually projected onto Watson quartic constants.

\section{What Has Been Verified So Far}
\subsection{1-mode benchmark}
The one-mode case is currently the main exact calibration benchmark.

The following facts have already been verified:
\begin{itemize}[leftmargin=*]
\item the harmonic block \(H_{12}H_{12}\) reproduces the expected second-order quartic structure,
\item the \(H_{22}\) block gives the expected \(B^2\)-type symbolic family in collapsed-coupling tests,
\item the full fourth-order one-mode result separates into the symbolic families
\[
A^2,\qquad AB C,\qquad A^2 C^2,\qquad B^2,
\]
where \(A\) denotes the collapsed \(\mu_1\) class, \(B\) the collapsed \(\mu_2\) class, and \(C\) the collapsed \(\Phi_3\) class,
\item in the one-mode benchmark, the purely cubic--cubic block \(H_{30}H_{30}\) can be imported exactly from the full BCH engine.
\end{itemize}

\subsection{2-mode benchmark}
The full BCH engine has also been run successfully for a collapsed-coupling two-mode case at fourth order. This already provides exact reference data for the \(H_{30}H_{30}\) symbolic class in a genuine multi-mode environment.

An important conclusion from the two-mode benchmark is:
\begin{quote}
the cubic--cubic multi-mode block is not simply a trivial scalar renormalization of the harmonic \(H_{12}H_{12}\) block.
\end{quote}
This means that a genuine multi-mode formula must be written for \(H_{30}H_{30}\); a one-factor rescaling is not sufficient.

\subsection{3-mode structure check}
A full exact three-mode fourth-order BCH run is still expensive, but the BCH engine has now been optimized enough to close targeted symbolic-class runs for the two crucial anharmonic channels \(H_{12}H_{30}\) and \(H_{30}H_{30}\). The main structural conclusion is:
\begin{quote}
for the cubic--cubic channel, the genuinely three-mode classes \(iij\_jkk\) and \(ijk\_ijk\) do appear and therefore must be part of the final multi-mode basis.
\end{quote}
This confirms that the reduced basis now implemented in the direct module is not over-parameterized: the extra three-mode invariants are physically real and not an artifact of the ansatz.

In addition, the targeted three-mode symbolic-class benchmarks are now cached and used directly by the channel solver.

\section{Current Code Status}
\subsection{Full engine}
The file
\texttt{derive\_watson\_quartic\_vanvleck.py}
currently contains:
\begin{itemize}[leftmargin=*]
\item the general operator algebra,
\item origin tracking of perturbative channels,
\item symbolic-family decomposition,
\item channel-aware pruning to reduce unnecessary intermediate growth,
\item targeted symbolic-class input selection for \texttt{--symbol-class-only} runs,
\item cached structural key-pair multiplication to accelerate the BCH commutator kernel.
\end{itemize}

This file remains the benchmark engine, but it is now also capable of targeted three-mode symbolic-class extractions for the channels that matter most in the direct architecture.

\subsection{Direct channel module}
The file
\texttt{quartic\_channels.py}
currently contains:
\begin{itemize}[leftmargin=*]
\item explicit direct formulas for \(H_{12}H_{12}\),
\item explicit direct formulas for \(H_{22}\),
\item a reduced calibrated basis for \(H_{12}H_{30}\), now fully fixed by one-, two-, and three-mode symbolic-class benchmarks,
\item a benchmark-backed minimal basis for \(H_{30}H_{30}\), obtained after testing reduced, full, scaffold, and symmetrized scaffold representations against an enlarged exact benchmark set,
\item explicit separation of the reduced invariants used for the mixed and cubic--cubic anharmonic channels,
\item explicit treatment of \(H_{40}\) as a structurally vanishing isolated quartic-force-field block in the quartic rotational sector,
\item exact one-mode and two-mode benchmark retrieval for the symbolic classes \(H_{12},H_{30}\), \(H_{30},H_{30}\), and \(H_{40}\),
\item exact three-mode benchmark retrieval for the symbolic classes \(H_{12},H_{30}\) and \(H_{30},H_{30}\),
\item a simple on-disk cache for expensive exact benchmark extractions.
\end{itemize}

More specifically, the current working reduced bases are:
\begin{itemize}[leftmargin=*]
\item for \(H_{30}H_{30}\), as an intermediate organizational level:
\[
\mathrm{iii\_iii},\qquad \mathrm{two\_mode\_mix},\qquad \mathrm{iij\_jkk},\qquad \mathrm{ijk\_ijk},
\]
with
\[
\mathrm{two\_mode\_mix}=\mathrm{iii\_iij}+\mathrm{iij\_iij}+\mathrm{iij\_ijj};
\]
\item for \(H_{12}H_{30}\):
\[
\mathrm{single\_mode},\qquad \mathrm{two\_mode\_mix},\qquad \mathrm{three\_mode\_mix}.
\]
\end{itemize}

For \(H_{12}H_{30}\), the one-mode benchmark fixes the diagonal part, the two-mode benchmark fixes the effective mixed-mode contribution, and the three-mode benchmark closes the reduced basis completely.

For \(H_{40}\), the present status is stronger than a numerical observation. If one considers the controlled limit
\[
H = H_{20}+H_{02}+H_{40},
\qquad
H_{12}=H_{22}=H_{30}=0,
\]
then the perturbation is purely vibrational and contains no rotational operators. Consequently, neither the BCH/Van Vleck commutator algebra nor the projected effective Hamiltonian can generate quartic rotational operators from \(H_{40}\) alone. In the present workflow this means
\[
\tau^{(H_{40})}=0
\]
for the quartic rotational block extracted from the effective Hamiltonian. Since the implemented projection is onto the vibrational ground-state block, this statement is used operationally for \(\langle 0|H_{\mathrm{eff}}|0\rangle\); however, the underlying reason is more general: a perturbation that is purely vibrational and independent of \(J\) cannot by itself create a \(J^4\) dependence.

For \(H_{12}H_{30}\), this refinement is now complete at the reduced-basis level: the coefficients
\[
\mathrm{single\_mode},\qquad \mathrm{two\_mode\_mix},\qquad \mathrm{three\_mode\_mix}
\]
are all fixed by the cached one-, two-, and three-mode symbolic-class benchmarks.

For \(H_{30}H_{30}\), the direct module was developed through a richer hierarchy of representations:
\begin{itemize}[leftmargin=*]
\item a reduced scalar mode basis,
\item a full six-class mode basis,
\item a rotationally resolved scaffold basis,
\item a symmetrized scaffold basis.
\end{itemize}

The first important conclusion is negative: the six-parameter symmetrized scaffold basis is not universal. Once an additional exact benchmark family at \(\mathrm{seed}=13\) is included, the corresponding fitted systems become inconsistent. Concretely, for each tensor component \(\tau_\mu\) the six-parameter system has
\[
\mathrm{rank}(M)=6,\qquad \mathrm{rank}(M|b)=7,
\]
so the problem is not an unresolved gauge freedom but an over-compression of the cubic--cubic basis itself.

The full 24-element rotational scaffold space remains consistent over the same benchmark set and has stable rank \(7\) for every \(\tau\) component. This identifies the correct intrinsic closure of the cubic--cubic channel: a seven-scaffold minimal basis. The stable pivot scaffolds are
\[
\mathrm{diag\_0\_iii\_iii},\quad
\mathrm{diag\_0\_iii\_iij\_1},\quad
\mathrm{diag\_0\_iii\_iij\_2},\quad
\mathrm{diag\_0\_iij\_iij\_1},\quad
\mathrm{diag\_0\_iij\_iij\_2},\quad
\mathrm{diag\_1\_iii\_iii},\quad
\mathrm{diag\_1\_iii\_iij\_0}.
\]

In other words, the apparent six-parameter gauge problem disappears once the benchmark space is enlarged: the correct statement is that \(H_{30}H_{30}\) requires seven independent scaffold structures at the \(\tau\)-tensor level. The earlier gauge-based narrative is therefore superseded by a stronger result: the cubic--cubic channel is now closed by a minimal seven-scaffold basis, not by a gauge-fixed six-scaffold one.

The corresponding explicit tensor equations are now recorded separately in the companion technical appendix
\texttt{tau\_channel\_equations.tex},
where the seven basis scaffolds are renamed in a more readable notation and all six components
\[
\tau^{(30,30)}_{xxxx},\quad
\tau^{(30,30)}_{yyyy},\quad
\tau^{(30,30)}_{zzzz},\quad
\tau^{(30,30)}_{xxyy},\quad
\tau^{(30,30)}_{xxzz},\quad
\tau^{(30,30)}_{yyzz}
\]
are written explicitly.

\section{Known Limitations}
At the current stage, the project is not yet complete.

The main limitations are:
\begin{itemize}[leftmargin=*]
\item the \(H_{12}H_{30}\) channel is operationally fixed in the reduced basis, but not yet derived as a closed analytic perturbative formula,
\item the \(H_{30}H_{30}\) channel is closed operationally in a benchmark-backed minimal seven-scaffold basis, but that basis has not yet been recast into a more compact analytic invariant structure derived from first principles,
\item the vanishing of isolated \(H_{40}\) in the quartic rotational block is theoretically motivated and consistent with the benchmark calculations, but a fully formal derivation in the same notation as the rest of the channel decomposition may still be worth recording separately,
\item the full BCH engine remains too expensive for routine three-mode fourth-order runs without substantial restriction.
\end{itemize}

\section{Immediate Next Step}
The next step is clear:
\begin{quote}
stabilize the direct-channel normal form by documenting the minimal seven-scaffold \(H_{30}H_{30}\) representation and then either
\begin{enumerate}[leftmargin=*]
\item derive a more compact analytic invariant basis equivalent to the seven-scaffold closure, or
\item accept the seven-scaffold basis as the operational working representation and move on to molecular input integration.
\end{enumerate}
\end{quote}

This is the best next move because:
\begin{itemize}[leftmargin=*]
\item the direct module is now complete enough to represent all channels observed to matter in the small exact benchmarks,
\item the remaining theoretical uncertainty is concentrated in finding a more transparent invariant reorganization of \(H_{30}H_{30}\), not in any unresolved numerical coefficient,
\item it avoids further expansion of the general BCH route except where it is genuinely needed for calibration.
\end{itemize}

\section{Long-Term Target}
The final intended workflow is:
\begin{enumerate}[leftmargin=*]
\item read molecular geometry, normal modes, and force constants from electronic-structure output,
\item compute \( \mu_{\alpha\beta,k} \) and \( \mu_{\alpha\beta,kl} \) analytically from inertia derivatives,
\item evaluate all quartic channels in a direct multi-mode implementation,
\item assemble the quartic tensor \( \tau \),
\item project the result to Watson quartic distortion constants.
\end{enumerate}

At that stage, the code will no longer depend on the expensive general BCH engine except as a benchmark and validation tool.

\section{Program Interface and Final Projection}
The channel-resolved VPT4 formulation developed in this project is intended to be interfaced with molecular data obtained from electronic-structure calculations and projected onto Watson quartic centrifugal distortion constants in a conventionally defined and program-independent manner.

\subsection{Required Molecular Data}
The program must read or reconstruct the following quantities:
\begin{enumerate}[leftmargin=*]
\item harmonic vibrational frequencies \(\omega_i\),
\item equilibrium geometry and atomic masses,
\item normal modes in a consistent mass-weighted convention,
\item cubic force constants \(\Phi_{ijk}\),
\item first derivatives of the inverse inertia tensor,
\[
\mu_{\alpha\beta,i}
=
\left.
\frac{\partial (I^{-1})_{\alpha\beta}}{\partial Q_i}
\right|_0,
\]
\item second derivatives of the inverse inertia tensor,
\[
\mu_{\alpha\beta,ij}
=
\left.
\frac{\partial^2 (I^{-1})_{\alpha\beta}}{\partial Q_i \partial Q_j}
\right|_0.
\]
\end{enumerate}

Within the present operational VPT4 framework, quartic force constants are not required as independent input for the quartic rotational tensor itself, because the isolated \(H_{40}\) channel does not generate quartic rotational operators.

\subsection{Channel Requirements}
The quartic rotational tensor is evaluated as
\[
\bm{\tau}
=
\bm{\tau}^{(12,12)}
+
\bm{\tau}^{(22)}
+
\bm{\tau}^{(12,30)}
+
\bm{\tau}^{(30,30)},
\]
with
\[
\bm{\tau}^{(40)} = \bm{0}.
\]
The required molecular data for each channel are:
\begin{itemize}[leftmargin=*]
\item \(H_{12}H_{12}\): \(\omega_i\) and \(\mu_{\alpha\beta,i}\),
\item \(H_{22}\): \(\omega_i\) and \(\mu_{\alpha\beta,ij}\),
\item \(H_{12}H_{30}\): \(\omega_i\), \(\mu_{\alpha\beta,i}\), \(\mu_{\alpha\beta,ij}\), and \(\Phi_{ijk}\),
\item \(H_{30}H_{30}\): \(\omega_i\), \(\mu_{\alpha\beta,i}\), and \(\Phi_{ijk}\).
\end{itemize}

\subsection{Projection Chain}
The final Watson reduction is not most naturally performed directly from the compressed six-component quartic tensor
\[
(\tau_{xxxx},\tau_{yyyy},\tau_{zzzz},\tau_{xxyy},\tau_{xxzz},\tau_{yyzz})
\]
alone. Instead, it is more naturally expressed through an intermediate Wilson-type quartic tensor.

In the present implementation the quartic operator block is first assembled as a full Wilson-type tensor,
\[
\tau_{ijkl},
\]
before any final compression. One then forms the reduced matrix
\[
\tau'_{ij},
\]
according to
\begin{align}
\tau'_{ii} &= \tau_{iiii},\\
\tau'_{ij} &= \tau_{iijj} + 2\tau_{ijij},
\qquad i\neq j.
\end{align}
The Cartesian quartic matrix is then
\[
T_{ij} = \tau'_{ij}/4.
\]

The cylindrical quartic quantities entering the Watson reduction are
\begin{align}
T_{400} &= \frac{3T_{11}+3T_{22}+2T_{12}}{8},\\
T_{220} &= T_{13}+T_{23}-2T_{400},\\
T_{040} &= T_{33}-T_{220}-T_{400},\\
T_{202} &= \frac{T_{11}-T_{22}}{4},\\
T_{022} &= \frac{T_{13}-T_{23}}{2}-T_{202},\\
T_{004} &= \frac{T_{11}+T_{22}-2T_{12}}{16}.
\end{align}

\subsection{Watson Reductions}
For asymmetric tops, the quartic Watson constants are
\begin{align}
\Delta_J   &= -T_{400}-2T_{004},\\
\Delta_K   &= -T_{040}-10T_{004},\\
\Delta_{JK}&= -T_{220}+12T_{004},\\
\delta_J   &= -T_{202},\\
\delta_K   &= -T_{022}-4\Sigma T_{004},
\end{align}
where \(\Sigma\) is the standard asymmetry-dependent reduction parameter. One first defines
\begin{equation}
\Sigma_1 = \frac{B-C}{2A-B-C},
\end{equation}
using the effective rotational constants \(A\), \(B\), and \(C\), and then for asymmetric tops
\begin{equation}
\Sigma = \frac{1}{\Sigma_1}.
\end{equation}

The symmetrically reduced constants are
\begin{align}
D_J   &= -T_{400}+\frac{1}{2}T_{022}\Sigma_1,\\
D_{JK}&= -T_{220}-3T_{022}\Sigma_1,\\
D_K   &= -T_{040}+\frac{5}{2}T_{022}\Sigma_1,\\
d_1   &= T_{202},\\
d_2   &= T_{004}+\frac{1}{4}T_{022}\Sigma_1,
\end{align}
where \(\Sigma_1\) is the inverse asymmetry parameter,
\[
\Sigma_1=\Sigma^{-1}=\frac{B-C}{2A-B-C}.
\]

\subsection{Practical Workflow}
A consistent implementation should therefore proceed as follows:
\begin{enumerate}[leftmargin=*]
\item read equilibrium geometry, atomic masses, harmonic frequencies, normal modes, and cubic force constants;
\item compute the first and second derivatives of the inverse inertia tensor from the equilibrium structure and the normal modes;
\item evaluate the perturbative channel contributions to the quartic rotational tensor;
\item assemble the full four-index quartic tensor \(\tau_{ijkl}\);
\item build the reduced matrix \(\tau'_{ij}\);
\item define \(T_{ij}=\tau'_{ij}/4\);
\item apply the transformation defined above to obtain Watson quartic centrifugal distortion constants.
\end{enumerate}

In this workflow the derivatives of the inverse inertia tensor are not assumed as direct input data. Instead, they are determined analytically from the molecular geometry, masses, and normal-mode displacements. This makes the procedure independent of any specific electronic-structure output format and suitable for implementation with data obtained from different quantum-chemical programs.

\section{Summary}
The project has moved from a purely general symbolic BCH implementation toward a more scalable and physically transparent channel-based architecture.

The present status can be summarized as follows:
\begin{itemize}[leftmargin=*]
\item the general symbolic engine exists and works as a benchmark tool,
\item the channel decomposition strategy is established,
\item one-, two-, and three-mode exact symbolic-class benchmarks are now available for the two key anharmonic channels \(H_{12}H_{30}\) and \(H_{30}H_{30}\),
\item \(H_{12}H_{30}\) is fixed operationally in the reduced basis,
\item \(H_{30}H_{30}\) is fixed operationally in a benchmark-backed minimal seven-scaffold basis,
\item the isolated \(H_{40}\) class is currently benchmarked as zero in the available small-case exact runs,
\item the next major technical task is to compress the seven-scaffold cubic--cubic closure into a more transparent analytic invariant form and then move to molecular input integration.
\end{itemize}

\end{document}
